Two-wheeled robots capable of balancing on a single wheel present a substantially harder control problem than conventional two-wheel balancing, with no closed-form model for the tilted equilibrium. This paper presents, to the best of our knowledge, the first reinforcement learning based controller for sustaining a two-wheeled robot in a stable single-wheel stance, trained end-to-end with PPO using a compact Gated Recurrent Unit (GRU) architecture that implicitly encodes system dynamics through its hidden state. The network is sized for real-time deployment on an STM32F446RE microcontroller at 50 Hz; domain randomization and a three-stage roll-perturbation curriculum bridge the sim-to-real gap. The learned policy proves capable of achieving and maintaining balance across a wide range of initial roll angles, reaching near-100\% survival probability within the trained $\pm 11°$ band around the 47° equilibrium stance.

%This paper investigates the use of deep reinforcement learning to solve the complex balancing problem of a wheeled-legged robot on a single wheel. This task represents a highly unstable, coupled, and underactuated control challenge, where pitch is managed via the drive wheel and roll is controlled through the leg joints. We utilize the Proximal Policy Optimization (PPO) algorithm within a physics simulation to train a robust control policy without requiring an explicit analytical model of the system's dynamics. Key features of our approach include a multi-phase learning curriculum to gradually broaden the range of starting conditions, and a reward function designed to let the agent autonomously discover its equilibrium roll angle rather than having it pre-defined. To bridge the Sim2Real gap, the solution incorporates domain randomization and is optimized for deployment on low-power embedded hardware, specifically an STM32F446RE microcontroller. Experimental results in simulation demonstrate a survival probability of 99.5\% at the optimal roll angle, confirming that the trained policy effectively masters the robot's inherent instability and successfully identifies its balance point.